\documentclass[11pt]{amsart}
\usepackage{geometry}                % See geometry.pdf to learn the layout options. There are lots.
\geometry{letterpaper}                   % ... or a4paper or a5paper or ... 
%\geometry{landscape}                % Activate for for rotated page geometry
%\usepackage[parfill]{parskip}    % Activate to begin paragraphs with an empty line rather than an indent
\usepackage{graphicx}
\usepackage{amssymb}
\usepackage{epstopdf}
\DeclareGraphicsRule{.tif}{png}{.png}{`convert #1 `dirname #1`/`basename #1 .tif`.png}

\title{Problema 2.1 }
%\date{}                                           % Activate to display a given date or no date

\begin{document}
\maketitle
%\section{}
%\subsection{}
\noindent El número de sonidos emitidos por un grillo en un minuto, es una función de la temperatura. Como resultado de esto, es posible determinar el nivel de la tempe­ ratura haciendo uso de un grillito como termómetro.\newline
La fórmula para la función es:\newline
\[
T = \frac{h}{4} + 40
\]
Donde: T representa la temperatura en grados Fahrenheit y N el número de sonidos emitidos por minuto
Construya un diagrama de flujo que le permita calcular la temperatura, te­ niendo en cuenta el número de sonidos emitidos por el grillo.\newline
Dato: N\newline

\noindent  \textbf{Explicación de las variables} \newline

\noindent N: Variable de tipo entero.\newline
T: Variable de tipo real. Almacena la temperatura en grados Fahren- heit.\newline
A continuación en la tabla 2.8 observamos el seguimiento del algoritmo para diferentes corridas.
\begin{table}[htbp]
\centering
\label{tab:corrida}
\begin{tabular}{|c|c|c|c|c|c|}
\hline
\multicolumn{3}{ |c| }{\textbf{Tabla 2.8}} \\ \hline
\textbf{Número de corrida} & \textbf{Dato - N} & \textbf{Resultado - T} \\ \hline
1 & 8& 42.00 \\ \hline
2 & 15& 43.75 \\ \hline
3 & 11& 42.75 \\ \hline
4 & 25 & 46.25 \\ \hline
5 & -5&  \\ \hline
6 & 50 & 52.50\\ \hline
\end{tabular}
\end{table}
\end{document}  